\section{Neo Labs 与 Robotics：新实验室、机器人和真实世界数据}

上一章讨论数据闭环，本章看硅谷新实验室和 Robotics。广密提到 SSI、Thinking Machines、Periodic Labs、Aethera、General Intuition、Reflection 等 Neo Labs，也讨论了机器人公司 Pi、Generalist、Sundee、1X、Figure、Google、Tesla 等。这一章的共同关键词仍然是数据：下一范式需要数据，机器人更需要真实世界数据。

\begin{figure}[H]
\centering
\includegraphics[width=0.88\textwidth]{neo-labs-map.png}
\caption{Neo Labs 分布：自主学习、Post-training、材料模型、多 Agent、世界模型数据和开源模型。自制概念图，依据 00:56:45--00:59:43 对谈内容整理。}
\end{figure}

\begin{knowledgebox}{读图：新实验室不是同一个方向的重复创业}
SSI 更像押自主学习；Thinking Machines 被视为 Post-training 与下一范式团队；Periodic Labs 聚焦材料模型；Aethera 探索 Multi-agent 环境；General Intuition 借助游戏高光数据做世界模型；Reflection 被描述为美国开源模型生态的补位。它们共同体现的是：OpenAI 人才溢出后，下一范式被多路线下注。
\end{knowledgebox}

\subsection{为什么 OpenAI 分裂后仍然强}

本节回到人才组织问题，因为 Neo Labs 的出现不是偶然的创业潮，而是前沿模型公司人才、路线和经验外溢的结果。理解这些团队从哪里来，才能判断它们为什么可能不只是“又一个模型公司”。

节目中有一个很有意思的观察：OpenAI 即使分裂出 Anthropic、Ilya 的 SSI、Mira/John Schulman 等相关团队，仍然很强。Anthropic 被视为 OpenAI 早期 Scaling team 的延伸，SSI 更接近 Pre-training/自主学习方向，Thinking Machines 更接近 ChatGPT 和 Post-training 经验。这说明 OpenAI 的人才密度和范式探索曾经足够集中，以至于外溢后形成多个新中心。

但这也意味着下一阶段不一定由原来的大公司独占。新实验室可能绕过现有模型追赶赛，直接押下一代学习方式、材料科学、多 Agent 环境或世界模型数据。这就是为什么本期既强调 OpenAI 领先，也强调不要只看现有榜单。

\subsection{Robotics：GPT-0.5 到 GPT-1 的阶段}

机器人部分的核心判断是：进步很快，但离 GPT-4 时刻还远。广密认为今天机器人可能处于 GPT-0.5 到 GPT-1，未来两三年可能到 GPT-2/GPT-3。原因在于机器人不像语言模型那样先统一再分化。语言模型有相对统一的 token、文本语料、Transformer 和预训练路径；机器人一开始就是分化的：不同硬件、不同传感器、不同场景、不同采集方式、不同评测目标。

机器人公司现在的分化体现在数据采集上。Google 和 Pi 有遥操作真机数据；Sundee 用手套和背包采家庭数据；Generalist 宣称有大规模真机数据；世界模型和 YouTube/游戏数据也开始进入训练。每条路径都在回答同一个问题：怎样让机器人获得足够多、足够真实、足够高质量的操作数据。

\begin{knowledgebox}{术语消化：机器人数据与泛化}
\begin{tabularx}{\textwidth}{p{0.22\textwidth}p{0.30\textwidth}X}
\toprule
术语 & 一句话解释 & 本期关系 \\
\midrule
遥操作 & 人远程控制机器人完成任务并记录过程 & 提供真实动作轨迹和环境反馈。 \\
真机数据 & 真实机器人在真实环境中的传感器和动作数据 & 比仿真更贴近落地，但采集昂贵。 \\
跨本体 & 不同机器人硬件之间共享或迁移能力 & 解决硬件不统一导致的数据碎片化。 \\
泛化 & 在未见过的新场景中仍能完成任务 & 机器人是否走出自动驾驶式长尾泥潭的关键。 \\
Scaling Law & 数据、模型规模和性能之间的可预测关系 & 如果机器人也出现规模律，产业节奏会更清楚。 \\
\bottomrule
\end{tabularx}
\end{knowledgebox}

\subsection{机器人为什么仍然卡在“场景”}

节目中提到，今天机器人主要可见场景仍集中在叠衣服、冲咖啡、家庭任务等演示或早期落地场景。硬件也变得越来越重要，有研究者甚至认为硬件可能占成功因素的很大比例。这提醒我们：机器人不是纯模型问题。它同时是感知、控制、硬件、成本、可靠性、安全、供应链、场景选择和售后运营问题。

\begin{warningbox}{机器人不是“把 LLM 接上手脚”}
语言模型的成功路径不能直接复制到机器人。机器人要处理物理摩擦、硬件差异、传感器噪声、安全风险、长尾场景和真实世界成本。没有泛化和持续学习，机器人会长期依赖人工采集和场景工程。
\end{warningbox}

\subsection{本章小结}

Neo Labs 和 Robotics 看似分散，其实都围绕下一阶段数据和学习效率展开。新实验室试图找到 Pre-training/RL 之外的新范式，机器人公司试图把真实世界操作变成可学习数据。二者的共同问题是：如何让模型从有限高质量数据中学到可泛化能力。

